\bookchapter{第三十章：投屏、演示与车载互联}

\chapterintro{手机屏幕再大也就7英寸，但你的电视有55英寸、你的投影仪有100英寸。把手机画面“扔”到大屏幕上，体验是完全不同维度的。这一章讲清楚投屏、演示和车载互联这三个高频场景，帮你把手机的能力延伸到更大的屏幕上。}

\booksection{一、无线投屏三大阵营——别搞混了}

无线投屏看起来都是“手机画面到电视上”，但底层技术有三种，互不通用：

\begin{center}
\begin{tabularx}{\linewidth}{|X|X|X|}
\hline
\textbf{技术} & \textbf{谁用} & \textbf{特点} \tabularnewline
\hline
AirPlay & iPhone/iPad/Mac & 苹果生态，体验最流畅。需要Apple TV或支持AirPlay的电视 \tabularnewline
Miracast & 安卓手机 & Wi-Fi Direct直连，不需要路由器。安卓手机"无线投屏"默认用这个 \tabularnewline
DLNA & 所有平台 & 只推送媒体文件（视频/音乐/图片），不镜像屏幕。适合推视频到电视上播放 \tabularnewline
\hline
\end{tabularx}
\end{center}

\booksubsection{各平台操作方法：}

\begin{itemize}
  \item \textbf{iPhone}：控制中心 > 屏幕镜像 > 选择你的电视/Apple TV。如果电视不支持AirPlay，可以买一个Apple TV盒子或支持AirPlay的电视棒。
  \item \textbf{安卓}：下拉通知栏 > 找到"无线投屏"/"智能投屏"/"多屏互动"按钮 > 选择电视。需要手机和电视连同一个Wi-Fi。
  \item \textbf{万能方案}：如果电视太老不支持任何无线投屏，买一个Chromecast或小米盒子/华为盒子插到电视HDMI口上，就能接收投屏了。
\end{itemize}

\booksection{二、有线投屏——更稳定、零延迟}

无线投屏方便但有延迟，玩游戏和做演示时体验不好。有线投屏能解决这个问题：

\begin{itemize}
  \item \textbf{iPhone}：需要一根Lightning/USB-C转HDMI的转换线（约100-200元）。插上后手机画面直接输出到电视/投影仪的HDMI口。
  \item \textbf{安卓}：部分手机支持USB-C转HDMI直连（需要手机支持DP Alt Mode，大部分旗舰都支持）。买一根USB-C to HDMI线即可。
  \item \textbf{会议室场景}：如果你经常需要用手机做演示，建议在包里常备一根转HDMI线。会议室的电视/投影仪几乎都有HDMI口，插上就能用，比折腾无线投屏靠谱多了。
\end{itemize}

\booksection{三、手机变电脑——桌面模式}

部分安卓手机支持“桌面模式”——连上显示器后，手机会输出一个类似电脑桌面的界面，可以用鼠标键盘操作：

\begin{itemize}
  \item \textbf{Samsung DeX}：三星旗舰的招牌功能。连接显示器后弹出桌面界面，可以多窗口操作，体验最接近真正的电脑。配合蓝牙键鼠，出差可以不带笔记本。
  \item \textbf{华为桌面模式}：华为Mate系列/Pura系列连接显示器后也有类似功能。支持多窗口、文件管理、办公App。
  \item \textbf{摩托罗拉Ready For}：Moto旗舰也支持，功能类似DeX。
  \item \textbf{适用场景}：轻办公（回邮件、改文档）、演示PPT、大屏看视频。不适合重度Excel操作或专业设计。
\end{itemize}

\booksection{四、车载互联——开车用手机的正确方式}

开车时操作手机是违法的，但导航、听音乐、接电话又是刚需。车载互联系统让你在车机屏幕上使用手机的核心功能，手机可以放在一边不碰：

\begin{itemize}
  \item \textbf{Apple CarPlay}：iPhone用户的首选。通过USB或无线连接车机，车机屏幕上出现iPhone的精简界面——导航（高德/百度）、音乐、电话、信息、播客。Siri语音控制全程免手。
  \item \textbf{Android Auto}：安卓用户对应方案。功能与CarPlay类似，支持Google Maps、Spotify等。国内因Google服务限制，部分功能不可用。
  \item \textbf{华为HiCar}：华为自研方案，国内车型支持率高。除了导航音乐，还能用华为视频通话、泊车影像融合等特色功能。
  \item \textbf{百度CarLife}：跨平台方案，iOS和安卓都能用。功能较少但兼容性广，很多老车型只支持CarLife。
\end{itemize}

\begin{tcbwarn}{安全提醒}
无论哪种车载互联方案，开车时都应尽量使用语音控制。导航提前设好目的地，音乐提前选好歌单，不要在行驶中操作车机屏幕。

\end{tcbwarn}

\booksection{五、会议室投屏实战——不再手忙脚乱}

开会是投屏最高频的场景，也是最容易出问题的场景。以下是几条实战经验：

\begin{itemize}
  \item \textbf{提前到场测试}：永远不要等到会议开始才连投屏。提前10分钟到场，测试好连接方式。Wi-Fi投屏失败的概率比你想的高得多。
  \item \textbf{准备有线方案做Plan B}：在包里放一根转HDMI线。无线投屏连不上的时候，有线就是救命稻草。
  \item \textbf{关闭通知预览}：投屏时你的微信消息会显示在会议室大屏上。投屏前在设置里关闭"锁屏通知内容预览"，或者开启"演示模式"（部分手机支持，投屏时自动隐藏通知）。
  \item \textbf{用"演示文档"而不是"操作手机"}：如果你要展示PPT，直接把PPT存到手机里用WPS打开全屏播放，而不是投屏整个手机桌面。这样更专业、更流畅。
\end{itemize}

\begin{tcbpractice}{动手：完成一次手机到电视的投屏}
\begin{enumerate}
  \item 确保手机和电视连同一个Wi-Fi。
  \item 用手机打开一张照片或一个视频。
  \item 尝试无线投屏（iPhone用屏幕镜像，安卓用无线投屏按钮）。
  \item 如果失败，检查电视是否支持对应协议，或尝试DLNA推送（在相册/视频App里点"投屏"按钮）。
  \item 投屏成功后，试试在手机上做其他操作，看看大屏是否同步显示。
\end{enumerate}

\end{tcbpractice}

\begin{tcbsummary}{本章小结}
\begin{itemize}
  \item 无线投屏三阵营：AirPlay（苹果）、Miracast（安卓）、DLNA（跨平台推媒体）
  \item 有线投屏更稳定零延迟，会议室场景必备转HDMI线
  \item Samsung DeX/华为桌面模式让手机变身轻办公电脑
  \item 车载互联：CarPlay/Android Auto/HiCar/CarLife四大方案
  \item 会议室投屏铁律：提前测试、准备Plan B、关闭通知预览
\end{itemize}

\end{tcbsummary}
